\NormalFont\ShortTitle{Tito}
{\MT Carta de Paulo a Tito

\par }\ChapOne{1}{\PP \VerseOne{1}Paulo, servo de Deus e apóstolo de Jesus Cristo, segundo a fé dos escolhidos de Deus, e o conhecimento da verdade, que é segundo a devoção divina;
\VS{2}em esperança da vida eterna, a qual Deus, que não pode mentir, prometeu antes dos princípios dos tempos,
\VS{3}E a seu
{\ADD{devido}} tempo
{\ADD{a}} manifestou: a sua palavra, por meio da pregação, que me foi confiada segundo o mandamento de Deus nosso Salvador.
\VS{4}Para Tito,
{\ADD{meu}} verdadeiro filho, segundo a fé em comum;
{\ADD{haja em ti}} graça, misericórdia
{\ADD{e}} paz de Deus Pai, e do Senhor Jesus Cristo, nosso Salvador.
\VS{5}Por esta causa eu te deixei em Creta, para que tu continuasses a pôr em ordem as coisas que estavam faltando, e de cidade em cidade constituísses presbíteros, conforme eu te mandei.
\VS{6}Se alguém for irrepreensível, marido de uma mulher, que tenha filhos fiéis, que não possam ser acusados de serem devassos ou desobedientes.
\VS{7}Porque o supervisor
\FTNT{*}{{\FR 1:7 }
supervisor = equiv. bispo
} deve ser irrepreensível, como administrador da casa de Deus, não arrogante, que não se ira facilmente, não beberrão, não agressivo, nem ganancioso.
\VS{8}Mas
{\ADD{que ele seja}} hospitaleiro, ame aquilo que é bom, moderado, justo, santo,
{\ADD{e}} tenha domínio próprio;
\VS{9}Retendo firme a fiel palavra que é conforme o que foi ensinado; para que ele seja capaz, tanto para exortar na sã doutrina, como
{\ADD{também}} para mostrar os erros dos que falam contra
{\ADD{ela}} .
\VS{10}Porque também há muitos insubordinados, que falam coisas vãs, e enganadores, especialmente aqueles da circuncisão;
\VS{11}Aos quais devem se calar, que transtornam casas inteiras, ensinando o que não se deve, por causa da ganância.
\VS{12}Um próprio profeta deles disse: Os cretenses sempre são mentirosos, animais malignos, ventres preguiçosos.
\VS{13}Este testemunho é verdadeiro; por isso repreende-os severamente, para que sejam sãos na fé.
\VS{14}Não dando atenção a mitos judaicos, e a mandamentos de homens, que desviam da verdade.
\VS{15}Realmente todas as coisas são puras para os puros; mas para os contaminados e infiéis, nada é puro; e até o entendimento e a consciência deles estão contaminados.
\VS{16}Eles declaram que conhecem a Deus, mas com as obras eles o negam, pois são abomináveis e desobedientes, e reprovados para toda boa obra.

\par }\Chap{2}{\PP \VerseOne{1}Tu, porém, fala o que convém à sã doutrina;
\VS{2}Aos velhos, que sejam sóbrios, respeitáveis, prudentes, sãos na fé, no amor e na paciência.
\VS{3}Às velhas, da mesma maneira,
{\ADD{tenham}} bons costumes, como convém a santas; não caluniadoras, não viciadas em muito vinho, mas sim instrutoras daquilo que é bom;
\VS{4}Para ensinarem às moças a serem prudentes, a amarem a seus maridos, a amarem a seus filhos;
\VS{5}A serem moderadas, puras, boas donas de casa, sujeitas a seus próprios maridos, para que a palavra de Deus não seja blasfemada.
\VS{6}Exorta semelhantemente aos rapazes, que sejam moderados.
\VS{7}Em tudo mostra a ti mesmo como exemplo de boas obras; na doutrina,
{\ADD{mostra}} incorrupção, dignidade, sinceridade;
\VS{8}Uma palavra sã
{\ADD{e}} irrepreensível, para que
{\ADD{qualquer}} opositor se envergonhe, nada tendo de mal para dizer contra vós.
\VS{9}Que os servos sejam sujeitos a seus próprios senhores, sendo agradáveis em tudo,
{\ADD{e}} não falando contra
{\ADD{eles}} .
\VS{10}Não
{\ADD{lhes}} furtando, mas sim mostrando toda a boa lealdade; para que em tudo adornem a doutrina de Deus nosso Salvador.
\VS{11}Porque a graça salvadora de Deus se manifestou a todos os homens.
\VS{12}Ensinando-nos que, ao renunciarmos à irreverência e aos maus desejos mundanos, vivamos neste tempo presente de maneira sóbria, justa e devota.
\VS{13}Aguardando a bem-aventurada esperança e o aparecimento da glória do grande Deus e nosso Salvador Jesus Cristo;
\VS{14}O qual deu a si mesmo por nós, para nos libertar de toda injustiça, e para purificar para si mesmo um povo particular, zeloso de boas obras.
\VS{15}Fala estas coisas, exorta, e repreende com toda autoridade. Ninguém te despreze.

\par }\Chap{3}{\PP \VerseOne{1}Relembra-os para se sujeitarem aos governantes e às autoridades, sejam obedientes,
{\ADD{e}} estejam preparados para toda boa obra.
\VS{2}Não insultem a ninguém, não sejam briguentos,
{\ADD{mas sim}} pacientes, mostrando toda mansidão para com todos os homens.
\VS{3}Porque nós também éramos tolos, desobedientes, enganados, servindo a vários maus desejos e prazeres, vivendo em malícia e inveja, detestáveis, odiando uns aos outros.
\VS{4}Mas quando a bondade e amor de Deus nosso Salvador para a humanidade apareceu,
\VS{5}Não pelas obras de justiça que nós tivéssemos feito, mas sim segundo sua misericórdia, ele nos salvou pelo banho do novo nascimento, e da renovação do Espírito Santo;
\VS{6}Ao qual ele derramou abundantemente em nós por meio de Jesus Cristo nosso Salvador;
\VS{7}Para que, ao termos sido justificados por sua graça, sejamos feitos herdeiros segundo a esperança da vida eterna.
\VS{8}{\ADD{Esta}} palavra
{\ADD{é}} fiel, e isto quero que insistas em confirmar, para que os que creem em Deus procurem se dedicar às boas obras; estas coisas são boas e proveitosas aos homens.
\VS{9}Mas evita as questões tolas, e às genealogias e discussões, e às disputas quanto à Lei, porque elas são inúteis e vás.
\VS{10}Ao homem rebelde, depois da primeira e
{\ADD{da}} segunda repreensão, rejeita
{\ADD{-o}} .
\VS{11}Sabendo que o tal está pervertido, e está pecando, estando a si mesmo condenado.
\VS{12}Quando eu enviar Artemas até ti, ou Tíquico, procura vir até mim a Nicópolis, porque eu decidi passar lá o inverno.
\VS{13}Auxilia com empenho a Zenas, o especialista na Lei, e a Apolo, na viagem deles, para que nada lhes falte.
\VS{14}E que também os nossos aprendam a se dedicarem às boas obras para os usos necessários, para que não sejam infrutíferos.
\VS{15}Todos os que estão comigo te saúdam. Saúda tu aos que nos amam na fé. A graça
{\ADD{seja}} com todos vós. Amém!
{\ADD{A carta a Tito, o primeiro bispo escolhido da igreja dos cretenses, foi escrita de Nicópolis na Macedônia}}
\par }